\section{Introduction}
\label{sec:intro}

Coarse particulate matter is a persistent air-quality problem across the Arabian Gulf. Mineral
dust from the surrounding deserts dominates the regional aerosol burden
\citep{gandham2020,shahid2026}, dust storms recur over the Emirates
\citep{barbulescu2020,karagulian2019}, Middle Eastern dust raises \pmten{} at urban receptors far
downwind \citep{givehchi2013,salmabadi2023}, and regional episodes have been reconstructed with
numerical models \citep{sulaymon2026}. Traffic, resuspended road dust, construction and industry
add an urban increment \citep{thorpe2008,amato2009,harrison2021,lenschow2001}; \pmten{} sources have
been apportioned for Kuwait \citep{hopke2026}, its variability described for Abu Dhabi
\citep{abuelgasim2020}, and deposition of particle-bound species shown to vary with meteorology and
surface geology \citep{bae2025be7}. The exposure matters for the burden of disease
\citep{cohen2017,lelieveled2015,chen2020,orellano2020,karanasiou2012}, yet monitors are sparse,
and where none exists assessments fall back on models \citep{shin2026resilient}: reanalyses such
as that of the Copernicus Atmosphere Monitoring Service (CAMS)
\citep{inness2019,flemming2015,morcrette2009} and MERRA-2 \citep{gelaro2017,randles2017}, or
fields that fuse models with satellite retrievals \citep{van_donkelaar2016,hammer2020}.
Meteorological reanalyses are routinely evaluated against stations, over Arabia as elsewhere
\citep{ullah2024,gomiscebolla2023}. For particulate matter this is harder, because aerosol
assimilation constrains the column optical depth, not the surface concentration. The optical
depth has been validated worldwide and over Arabia \citep{eskes2015,gueymard2020,labban2023}, but
surface particulate matter in the Middle East was found to differ substantially \citep{ukhov2020}.
Whether such a difference lies in the timing of dust events or in the surface concentration
attributed to them cannot be decided without an independent record of dust.

Two further issues decide what an evaluation against monitors can conclude. The first is the
averaging interval of the observations. A model value and an observation are comparable only if
they are averaged over the same interval \citep{schutgens2017}, and open archives that
redistribute regulatory data \citep{rosales2025} do not always state the interval over which a
network reports; where an air-quality index is based on a running mean, the archived value may be
that mean rather than an hourly concentration. The second is representativeness. A grid value is
an average over several thousand square kilometres and a monitor samples one location inside
that area, so the two would differ even if the model were perfect
\citep{schutgens2016,schutgens2017,janjic2018}; the area that a monitor represents depends on the
sources around it \citep{janssen2012,righini2014,duyzer2015}, and the representativeness of
column aerosol measurements has itself been mapped \citep{chen2025aod}. For coarse particles in
an arid city the difference can be large, because roads, construction sites, industrial yards
and bare soil are much smaller than a grid cell and the particles they emit deposit within
kilometres. The effect of sub-grid variability on exposure estimates has been quantified for
Europe \citep{denby2011,denby2024} and addressed by sub-grid downscaling of regional models
\citep{theobald2016}, and local and long-range contributions have been separated with dense
sensor networks \citep{defoy2024}, but over the Gulf the sub-grid variability of \pmten{} has
rarely been measured, because it takes several monitors inside one model cell to measure it.

The usual remedy for model error is statistical. A correction model is fitted to the difference
between observations and the gridded field at monitored locations and is then applied wherever
the field is available \citep{djalalova2015,huang2017}, and machine-learning models of air
quality raise issues that are easily overlooked \citep{tang2024ml}; for CAMS this has been done for
\pmten{} in dust forecasting \citep{sarafian2025} and for ozone at urban stations
\citep{kaffashzadeh2024}. Alternatively, point measurements are interpolated into continuous
fields, whose accuracy depends on the spatial structure of the field being mapped
\citep{kang2026mapping}. Whether such a correction
carries to places without monitors is normally judged by withholding one station at a time. In a
dense network that test is weaker than it looks, because the withheld station's neighbours remain
in the training set a few kilometres away, sampled on the same days and often inside the same
model cell \citep{roberts2017,ploton2020,meyer2021aoa,wadoux2021}, a problem for which
dedicated designs of spatial cross-validation have been proposed
\citep{valavi2019,mila2022,linnenbrink2024,meyer2022global}. The question is the
representativeness question seen from the other side: if what the gridded field misses is local,
a correction can only be learned from monitors that share that locality.

The regulatory network of the Emirate of Abu Dhabi allows these questions to be put to data.
Nineteen \pmten{} monitors lie at a median nearest-neighbour separation of 6.5~km, and sixteen of
them share a $0.75^{\circ}$ CAMS cell with at least one other monitor. We use this network to ask
(i) over what interval the archived values are averaged, and what follows for comparison with a
model; (ii) whether the CAMS reanalysis reproduces the timing of dust and its surface
concentration, judged against an independent sun photometer and airport weather reports;
(iii) which part of the observed variability the reanalysis reproduces, separating the
day-to-day level of the region from the differences between places, and how much of the spatial
variability lies inside single model cells; (iv) whether a statistical correction fitted at
monitored sites improves the reanalysis at a site treated as unmonitored; and (v) what the
answers imply for assessing exceedance of the World Health Organization 24-h guideline\footnote{World Health Organization, 2021. WHO global air quality guidelines: particulate matter (PM$_{2.5}$ and PM$_{10}$), ozone, nitrogen dioxide, sulfur dioxide and carbon monoxide. ISBN 978-92-4-003422-8. \url{https://www.who.int/publications/i/item/9789240034228}.}. A recent evaluation showed
that global reanalyses cannot reproduce the seasonal and diurnal cycles of surface ozone in a
data-sparse tropical region \citep{vieira2023}; the present study addresses the spatial
counterpart of that problem for coarse particles in an arid one.

\section{Data and methods}

\subsection{Surface observations}
\label{sec:obs}

\pmten{} from the 19 monitors operated by the Environment Agency -- Abu Dhabi was taken from the
OpenAQ open-data archive \citep{rosales2025} (Figure~\ref{fig:map}, Table~\ref{tab:stations}).
The archive holds 98,459 hourly \pmten{} values for these monitors between 29~October 2022 and
20~May 2024. Its coverage is uneven: from October 2022 to January 2023 it holds 14 of the 24
hours of the day, in February, May and July 2023 a single value per day at 08~UTC, in April and
June 2023 twelve hours, and from August 2023 onwards all hours; March 2023 is absent
(Supplementary Information, Section~S1). No value is negative and no upper cap was applied.

Two observation records are formed from the archive. The \emph{three-hourly record} holds the
archived value at every three-hourly time step of the reanalysis, 32,566 values, between 1,678
and 1,774 per monitor. The \emph{daily record} holds one value per monitor and day, the value at
08~UTC, which is the only hour the archive holds in every month: 5,989 monitor-days on 331 days.
The 20,818 three-hourly values up to December 2023 that were analysed in earlier work on this
network were checked against the archive and are identical to the archived values at the same
time stamps.

Each monitor carries the type assigned by the network operator and published by the Statistics
Centre -- Abu Dhabi\footnote{Statistics Centre -- Abu Dhabi, Air Quality Statistics: Methodology. \url{https://scad.gov.ae} (accessed 13 November 2026).}: industrial (four monitors), traffic (three), urban
(three), suburban (seven) and rural background (two). We use this classification as given and do
not reassign any station. For 2017--2018 the network is reported to have measured \pmten{} with
OPSIS SM200 monitors \citep{abuelgasim2020}, which determine particle mass by beta
attenuation\footnote{Battelle, Environmental Technology Verification Report: Opsis AB SM 200
Automatic Particle Monitor, U.S. Environmental Protection Agency, 2001.
\url{https://archive.epa.gov/research/nrmrl/archive-etv/web/pdf/01_vr_opsis_sm200.pdf}
(accessed 14 November 2026).}. The archive does not record the instrument, the inlet height or
changes of equipment for the period studied.

\subsection{The archived values are 24-h running means}
\label{sec:running}

The archived \pmten{} values are not hourly concentrations. Four properties of the record show
that each is an average over about 24~h (Section~S2). (i) Successive hourly values are almost
identical: the lag-1~h autocorrelation has a median of 0.996 across the monitors, all values are
integers, and 35\% of them equal the value one hour earlier. Nitrogen dioxide, sulfur dioxide,
ozone and carbon monoxide reported by the same stations in the same files have lag-1~h
autocorrelations of 0.80 to 0.91. (ii) For a trailing 24-h running mean $x$ of hourly
concentrations $h$, $x_t-x_{t-1}=(h_t-h_{t-24})/24$, so each hourly increment is undone 24~h
later and the increments are negatively correlated with themselves at that lag. The
autocorrelation of the hourly \pmten{} increments at 24~h is negative at all 19 monitors (median
$-0.36$, range $-0.53$ to $-0.19$); for the four gases the median is positive ($+0.06$ to
$+0.33$), as a diurnal cycle requires. (iii) The observed values correlate better with the mean
of the reanalysis over the preceding 24~h than with the reanalysis at the same instant (median
Pearson $r$ over the monitors 0.50 against 0.37), and in the three-hourly record to December 2023
the correlation falls at 17 of 19 monitors as the 24-h window is moved from trailing to centred
to leading. (iv) The national guideline for the Emirati Air Quality Index prescribes a 24-h
rolling average for \pmten{}, formed from at least 18 valid hourly values\footnote{Guideline for calculating the Emirati Air Quality Index (EAQI), 2023. \url{https://airquality.ncm.gov.ae/resources/pdf/aqi-quickguide-en-2023.pdf} (accessed 13 November 2026).}.

We therefore treat every archived value as the mean concentration over the 24~h ending at its
time stamp. Three consequences follow. The model counterpart of an archived value must be averaged
over the same 24~h (Section~\ref{sec:cams}). Values three hours apart share 21 of their 24 hours,
so that the daily record is used wherever independence in time matters. And the record contains
no information on the diurnal cycle or on peaks shorter than a day; the hourly concentrations
cannot be recovered from integer-rounded running means.

\subsection{Reanalysis and matched averaging}
\label{sec:cams}

Surface \pmten{} from the CAMS global reanalysis EAC4 \citep{inness2019}, whose aerosol scheme is
described by R\'emy et al.~\citep{remy2019}, was retrieved from the Copernicus Atmosphere Data Store at
$0.75^{\circ}$ and three-hourly resolution for 1~October 2022 to 31~May 2024 without gaps, and
converted to \ugm{}. Each monitor was matched to the nearest grid point; bilinear interpolation
between the four surrounding grid points was used as a check (Section~S6). The 19 monitors occupy
nine cells; six cells hold two or more monitors and account for 16 of the 19
(Figure~\ref{fig:map}). Fifteen of the 171 monitor pairs share a cell, at separations from 1.0 to
46~km. The retrieved series reproduces, to within $2\times10^{-5}$~\ugm{}, all 20,818 model values
used in earlier work on this network.

For an archived value at time $t$ the model counterpart is the mean of the eight three-hourly
model values in the 24~h ending at $t$. Because the model series is continuous, every value of
both observation records has a counterpart. Comparing the archived values with the model at the
same instant instead, as an hourly record would be treated, changes the mean bias of a monitor by
at most 2~\ugm{}, but lowers the correlation in time and, as Section~\ref{sec:res-transfer}
shows, alters the verdict on statistical correction.

\subsection{Error structure, sub-grid variability and the regional level}

For each monitor we computed, on the three-hourly record, the mean bias (CAMS minus observed),
mean absolute error, root-mean-square error, and the Pearson and Spearman correlations between
the matched model and observed series. Spatial organisation of the mean bias was measured with
Moran's $I$ \citep{moran1950} in distance bands, with significance from 9,999 permutations of the station labels.

Sub-grid variability was measured in three ways. First, for each cell holding two or more
monitors we took the range and standard deviation of the monitors' mean observed concentrations,
and partitioned the between-monitor sum of squares of mean concentration over all 19 monitors
into a within-cell and a between-cell part; the partition was repeated on station medians, which
a single episode cannot dominate. Second, for every pair of monitors we computed the Spearman correlation \citep{spearman1904} of their observed three-hourly series over common time steps (at least 100) and
compared same-cell pairs with the correlation between each monitor and its own matched model
series. Third, for each pair we counted how often the two monitors lay on opposite sides of the
guideline value of 45~\ugm{}, over common time steps and over days of the daily record. A gridded
product that assigns the two monitors one value is necessarily wrong for one of them on every such
occasion.

To separate the timing of high concentrations from their location, each day of the daily record
with values from at least ten monitors was characterised by the median over the monitors (the
\emph{regional level}), for the observations and for the matched model values. Days were divided
into terciles of the observed regional level, and the bias and the spread between monitors were
computed within each tercile. Because the regional level has a strong annual cycle, the terciles
were also formed separately within the warm season (May to September) and the cool season
(November to February). The difference between the bias of the metropolitan and of the Al~Ain
monitors was computed with four aggregations (mean over monitor-days, mean and median over
per-monitor means, and without the most biased monitor), and its dependence on the regional level
was estimated by regressing its daily value on the regional level (days with at least four
metropolitan and two Al~Ain monitors; 10,000 bootstrap resamples of days).

Dust was identified independently of the \pmten{} record with the Aerosol Robotic Network
(AERONET) of sun photometers \citep{holben1998}. Of the 24 AERONET sites in the Emirates, only DEWA
Research Centre (24.77$^{\circ}$N, 55.37$^{\circ}$E, in Dubai, 106~km from central Abu Dhabi and
71~km from Al~Ain) holds data through the study period. We used its daily averages of the
coarse-mode aerosol optical depth at 500~nm from the spectral deconvolution algorithm
\citep{oneill2003}, Version~3, Level~1.5 \citep{giles2019}, averaged over the day of the
observation and the previous one so as to cover the 24-h window, and divided the days that have
such a value into terciles of it. Inside the network, the hourly weather reports (METAR) of Abu
Dhabi International Airport, 7.5 to 30~km from the seven metropolitan monitors, were used in the same
24-h windows: a day counts as a dust day when at least one report carries a dust or sand code, and
as clear when no report without precipitation, fog or mist has a visibility below 10~km. From the reanalysis itself we took the total, dust and
sea-salt aerosol optical depth at 550~nm and the 10-m wind, averaged over the same 24~h and over
the three cells that hold the metropolitan monitors, and related the daily metropolitan minus
Al~Ain difference to the dust optical depth, the wind speed and the wind sector. In the cell of the sun
photometer the total and dust optical depth of the reanalysis were also compared in magnitude with
the total and coarse-mode optical depth measured there, and the regional level per unit of column
dust was compared between model and observations. The analyses by
season, aggregation, interpolation, sun photometer, airport reports, reanalysis composition and
wind, and column magnitudes were added after the main results were known; they are exploratory, and their $p$-values are
not corrected for multiplicity. The stability of the results over time was checked by repeating the daily comparison for four sub-periods.

The association between station type and bias was tested with a Kruskal--Wallis test across the
five types \citep{kruskal1952} and, separately, with a one-sided Mann--Whitney test
\citep{mann1947} and a 200,000-draw permutation
test contrasting the monitors typed industrial, traffic or urban with those typed suburban or
rural background. The two-group contrast was formed after the five-group result was known; both
are reported.

\subsection{Statistical correction at a monitor treated as unmonitored}
\label{sec:methods-transfer}

The correction is a gradient-boosted regression \citep{friedman2001,pedregosa2011} of the
observed-minus-model difference on the matched model value, the sine and cosine of the hour and
of the day of the year, and latitude and longitude; a coordinate-free variant omits the last two.
Each monitor in turn is treated as unmonitored: the regression is trained on other monitors only
and its prediction, added to the model value, is compared with the withheld observations. The
gain is expressed as $\Delta=\log(\mathrm{MAE_{raw}}/\mathrm{MAE_{corrected}})$, so that
$\Delta>0$ means the correction reduced the mean absolute error. The three-hourly record is used,
after a screen that removes runs of six or more identical consecutive values (431 values, 1.3\%),
which leaves 32,135 values.

Withholding a monitor does not by itself make it unmonitored, because information about it
remains in the training data. Three kinds are removed in turn:
\begin{enumerate}
\item \emph{Nearby monitors.} All training monitors within a radius $R$ of the target are
      withheld, for $R=0$, 25, 50 and 100~km; $R=0$ is leave-one-station-out.
\item \emph{Observations from the same days.} The target record is cut into six blocks holding
      equal numbers of values, and each block is predicted from training data lying outside the
      block and outside a guard interval of ten days on either side. The matched residual stays
      correlated at this lag (median autocorrelation $+0.37$ at one day, $+0.18$ at three days
      and $+0.12$ at ten days), so no guard within the record makes the training data independent
      of the block; a three-day guard is reported as a check (Section~S5).
\item \emph{Monitors sharing the model cell.} Every training monitor in the target's CAMS cell
      is withheld, because it carries the same model series as the target.
\end{enumerate}
The test is an exact two-sided sign test on the 19 per-monitor values of $\Delta$, corrected by
the Benjamini--Hochberg procedure \citep{benjamini1995} across the 24 combinations of radius,
feature set and stage. Intervals for the median $\Delta$ are from 2,000 bootstrap resamples of
the nine occupied cells, the unit within which monitors share a model series \citep{efron1979}.
To show what the averaging interval does to this analysis, it was repeated on the same values
with the instantaneous model value in place of the matched one.

\subsection{Exceedance of the 24-h guideline}

Because the archived values are 24-h means, exceedance of the WHO 24-h guideline value can be
read from them directly. The daily record was used, with its matched model counterpart: 5,989
monitor-days. The observed value was classed as exceeding 45~\ugm{} or not and compared with the
same classification of the matched reanalysis. We report the probability of detection (POD), the
probability of false detection (POFD) and two scores that discount agreement expected by chance,
the equitable threat score (ETS) and the Heidke skill score (HSS); both are equitable only in
the limit of large samples \citep{hogan2010}. Intervals are from 2,000
bootstrap resamples of monitors. Two references that use the monitor itself indicate what local
information adds: the reanalysis shifted by the mean observed-minus-model difference at that
monitor over its preceding eight daily values, and persistence of the monitor's value from the
previous day. Neither is available where there is no monitor.

\section{Results}

\subsection{The reanalysis places dust in time but falls short of its surface concentration}
\label{sec:res-dust}

The reanalysis follows the day-to-day course of \pmten{} over the region. The observed regional
level and its model counterpart are related from day to day (Spearman $\rho=0.67$ over 327~days;
Figure~\ref{fig:regional}a), and the relation holds in each sub-period ($\rho=0.55$ to 0.71). At
single monitors the median Spearman correlation between the matched model and observed series is
0.55 (range 0.26 to 0.68; median Pearson $r=0.50$). This lies between the correlation of two
monitors that share a cell (median 0.66 over the 15 pairs) and that of monitors in different cells
(0.48; Figure~\ref{fig:subgrid}b): in time, the model stands in for a monitor somewhat better than
a distant monitor does, and somewhat worse than a neighbour.

The days on which the regional level is high are days of dust. Terciles of the observed regional
level are largely seasonal: 77 of the 107 days of the highest tercile fall between May and
September, and 94 of the 109 days of the lowest between November and February. They nonetheless
follow dust. On the 228 days with a sun-photometer value, the coarse-mode optical depth correlates
with the observed regional level at $\rho=0.70$ (with the modelled one at 0.54), and of the 71 days
of the highest regional tercile among them, 51 fall in the highest tercile of coarse-mode optical
depth and none in the lowest (Section~S7). The airport reports inside the network agree. On the
14 days with a reported dust event the regional level has a median of 125~\ugm{}, against 69 on
the other days, and the model falls short of it by 37~\ugm{}, against 21 (Section~S7).

The reanalysis times these events correctly in the column. Its dust optical depth in the cell of
the sun photometer follows the coarse-mode optical depth measured there (Spearman $\rho=0.85$ over
226 days), and its total optical depth matches the measured one in magnitude (median ratio 1.07;
Section~S8). At the surface it falls increasingly short. Classed by the coarse-mode optical depth,
the median difference between the modelled and the observed regional level is $-14$, $-26$ and
$-32$~\ugm{} from the least to the most dusty third of days. The model's day-to-day range is
compressed (regression slope of the model network median on the observed one 0.55), and its network
median exceeds 150~\ugm{} on 2 of the 327 days, against 12 in the observations. Two steps between
column and surface contribute (Section~S8). The reanalysis assigns dust a smaller part of the
column than the sun photometer assigns the coarse mode (median 0.31 against 0.48), so that its dust
optical depth is lower than the coarse-mode optical depth on 73\% of days (median ratio 0.76); and
per unit of column dust it carries less regional \pmten{} at the surface than is observed (median
ratio 0.86, Wilcoxon $p=0.03$).

\subsection{The bias is organised in space, and the conurbation carries a local deficit}
\label{sec:res-bias}

Mean observed \pmten{} ranges from 54 to 179~\ugm{} across the network
(Table~\ref{tab:stations}). Every one of the 19 station means lies above the 24-h guideline value
of 45~\ugm{}, and between 58\% and 99\% of the archived 24-h means at a monitor exceed it. Of the
5,989 monitor-days, 84\% exceed the guideline, from 60--68\% in the cleanest winter months to
more than 90\% from May to September 2023 (Section~S1), in line with the summer maximum reported
for the emirate in 2017--2018 \citep{abuelgasim2020}.

The matched reanalysis places the same locations between 37 and 106~\ugm{}. Its mean bias ranges
from $-128$~\ugm{} at Al~Mafraq to $+40$~\ugm{} at Sweihan, with a station mean of $-27$~\ugm{};
the median absolute error is 30~\ugm{} and the median root-mean-square error 48~\ugm{}. The model
does not reproduce which locations are high. Across the 19 monitors the standard deviation of
mean concentration is 27~\ugm{} in the observations and 19~\ugm{} in the model, and the two sets
of means are, if anything, inversely related (Spearman $\rho=-0.45$, $p=0.05$): the model's
highest values lie over the eastern desert, the observed ones in and around the city of Abu Dhabi.
None of this depends on taking the nearest grid point. In the daily record, bilinear interpolation
changes the mean bias of a monitor by at most 21~\ugm{}, leaves the ranking of the monitors by
bias almost unchanged (Spearman 0.95) and the station means of observations and model still
inversely related ($\rho=-0.43$), and narrows the spread of the model between monitors rather
than widening it (Section~S6).

The bias is not scattered at random. Moran's $I$ of the station mean bias is 0.80 for pairs
within 15~km ($p=0.001$) and 0.41 for pairs within 100~km ($p<0.001$). The seven monitors of the
Abu Dhabi metropolitan area are all underestimated, by 38 to 128~\ugm{}, whereas the four
monitors in Al~Ain, more than 100~km inland, are reproduced to within $-7$ to $+5$~\ugm{}
(Figure~\ref{fig:map}). The operator's station typology accounts for part of this pattern but not
cleanly. Monitors typed industrial, traffic or urban are underestimated by a median 55~\ugm{} and
suburban and rural-background monitors by a median 4~\ugm{} (Mann--Whitney $p=0.019$; permutation
$p=0.015$; Cohen's $d=-1.04$), but the five types taken separately do not differ significantly
(Kruskal--Wallis $p=0.23$), and each source-oriented type contains a monitor with little bias:
Ruwais (industrial, $-1$~\ugm{}), Al~Ain Street (traffic, $-7$~\ugm{}) and Zakher (urban,
$+5$~\ugm{}). Type alone does not predict where the reanalysis fails; location does. The pattern
is stable in time: the ranking of the monitors by daily mean bias in each of four sub-periods
agrees with that of the whole record (Spearman 0.92 to 0.98; Section~S4).

On top of the regional deficit of Section~\ref{sec:res-dust}, which the coast shares with the
interior, the conurbation carries a deficit of its own. When the days are divided by observed
regional level, the Abu Dhabi metropolitan monitors are underestimated by 46~\ugm{} in the lowest
tercile (regional level 24 to 61~\ugm{}), by 58~\ugm{} in the middle one and by 86~\ugm{} in the
highest (87 to 257~\ugm{}), while the Al~Ain monitors move from $+13$ through 0 to $-23$~\ugm{}
(Figure~\ref{fig:regional}b). The pattern does not arise from the season: with terciles formed
within each season, the metropolitan bias moves from $-39$ to $-63$~\ugm{} in the cool season and
from $-50$ to $-115$~\ugm{} in the warm one, and the Al~Ain bias from $+15$ to $-12$ and from $+2$
to $-41$~\ugm{} (Section~S6). The difference between the metropolitan and the Al~Ain bias is
$-59$, $-57$ and $-63$~\ugm{} in the three terciles. On the cleanest third of days this local
deficit lies between $-50$ and $-62$~\ugm{} however it is aggregated. On the highest days its size
depends on the aggregation (Section~S6): $-41$~\ugm{} without Al~Mafraq, the monitor with the
largest bias, $-33$~\ugm{} with the median over monitors, and $-75$~\ugm{} in the highest tercile
of the warm season. Regressed day by day on the regional level, it has a slope of $-0.00$ (95\%
interval $-0.23$ to $+0.22$) over all monitors and $+0.18$ ($-0.02$ to $+0.39$) without
Al~Mafraq. Whether the local deficit changes with the regional level is therefore not
established; that it is present on clean days is.

The local deficit does not follow dust. Classed by the coarse-mode optical depth, the metropolitan
minus Al~Ain difference is $-60$, $-58$ and $-53$~\ugm{} ($-53$, $-46$ and $-30$~\ugm{} without
Al~Mafraq) from the least to the most dusty third of days, and on the least dusty third every
metropolitan monitor is underestimated, by 30 to 97~\ugm{} (Section~S7). On the days with a
reported dust event the metropolitan minus Al~Ain difference is $-72$ against $-58$~\ugm{}
($p=0.26$). On the 150 clear days the difference is $-60$~\ugm{} and every metropolitan monitor
is underestimated, by 33 to 137~\ugm{} (Section~S7). Nor does it follow the reanalysis's own
fields (Section~S8). The metropolitan minus Al~Ain difference is unrelated to the dust optical
depth over the metropolitan cells ($\rho=-0.01$; $+0.21$ without Al~Mafraq, the deficit becoming
smaller on dustier days) and to the 10-m wind speed ($\rho=+0.06$; $+0.13$ without Al~Mafraq), and
it is the same with wind from the north-west quadrant, which prevails on 224 of 324 days, as with
wind from elsewhere (medians $-58$ and $-60$~\ugm{}, Mann--Whitney $p=0.22$). Sea salt accounts
for a median 1.5\% of the modelled optical depth over the metropolitan cells.

Across the regional terciles the spread between monitors grows in the observations but hardly
in the model: the day-by-day standard deviation across the network has a median of 27, 31 and
61~\ugm{} in the three terciles in the observations, against 21, 22 and 25~\ugm{} in the model.

\subsection{Up to half of the spatial variance lies inside single model cells}
\label{sec:res-subgrid}

Table~\ref{tab:cells} and Figure~\ref{fig:subgrid}a show the six cells that hold more than one
monitor. In four of them the station means agree to within 1--16~\ugm{}. In the other two they do
not: station means differ by 56~\ugm{} in the cell containing E11~Road, Habshan South and Bida
Zayed, and by 89~\ugm{} in the cell containing Al~Mafraq, Mussafah and Baniyas, with within-cell
standard deviations of 30 and 46~\ugm{}. Over the whole network 46\% of the between-monitor
variance of mean \pmten{} lies within cells. The two high-spread cells differ in kind. The spread
in the E11~Road cell comes mainly from one episode between late August and mid-September 2023,
during which the 24-h mean at that roadside monitor reached 2,978~\ugm{}; its median (61~\ugm{})
is close to those of its neighbours, and on station medians the within-cell share of the variance
is 31\%. The spread in the Al~Mafraq cell, an industrial area south-east of Abu Dhabi island, is
persistent: the station medians there differ by more than 70~\ugm{}, and the range of station
means exceeds 100~\ugm{} in each sub-period of the daily record. In every sub-period between 42\%
and 51\% of the variance of station means, and between 35\% and 44\% of that of station medians,
lies within cells (Section~S4). No field on this grid, however accurate as a cell average, can
represent that part of the variability. Proximity alone does not guarantee coherence either: the two monitors
closest to each other, Khadeeja School and Hamdan Street, 1.0~km apart with mean concentrations of
95 and 93~\ugm{}, correlate at only 0.33.

The six cells also separate two kinds of error that a single bias statistic merges. In the two
cells over Abu Dhabi island and its mainland suburbs the paired monitors agree closely (ranges of
1~\ugm{} each), yet the model lies 56 and 55~\ugm{} below them. There the discrepancy is not the
idiosyncrasy of one monitor; it is an urban-scale enhancement that the cell value does not carry.
In the two high-spread cells, by contrast, the monitor closest to the model is within 6 and
38~\ugm{} of it and the largest errors belong to individual industrial and roadside sites. In
Al~Ain the four monitors agree with each other and with the model.

For the guideline the consequence is direct. Two monitors in the same cell lie on opposite sides
of 45~\ugm{} in a median 10\% of common values, and in up to 39\% for the most dissimilar pair;
over days the figure is 14\% of 4,590 pair-days. A product with one value per cell is necessarily
wrong about exceedance at one of the two monitors on those days.

\subsection{A statistical correction gives no established gain, and the averaging interval
decides whether one appears}
\label{sec:res-transfer}

With each monitor withheld in turn and nothing else removed, the correction reduces the error at
13 of 19 monitors with coordinates as predictors (median $\Delta=+0.19$, cell-bootstrap interval
$+0.04$ to $+0.32$) and at 12 of 19 without them ($+0.11$, $-0.03$ to $+0.42$;
Figure~\ref{fig:transfer}a, Table~\ref{tab:transfer}). The sign test establishes neither
($p=0.17$ and 0.36 before correction for multiplicity). Withholding the monitors within 25~km
already brings the median to $+0.00$ for both feature sets. When the target record is predicted
in time blocks the gain disappears even with every neighbour retained (medians $+0.04$ and
$+0.02$, 11 and 10 of 19 monitors improved), and withholding the monitors of the target's cell
changes little further. None of the 24 comparisons differs from zero ($q=1.00$ for each). A
three-day guard gives the same result, with medians between $-0.28$ and $+0.09$ in the stages
with time blocks and no comparison different from zero (Section~S5).

The averaging interval is decisive for this conclusion. If the same 32,135 archived values are
compared with the model at the same instant, as an hourly record would be treated, the
leave-one-station-out gain is $+0.32$ with coordinates and $+0.42$ without, an apparent reduction
of the error by 28\% and 34\% at 17 and 16 of 19 monitors ($q=0.009$ and 0.015). This apparent
gain survives the controls that remove it under matched averaging: 16 of the 24 comparisons
establish it, including the coordinate-free correction with time blocks, the target's cell and
all monitors within 100~km withheld ($+0.15$, 15 of 19 monitors, $q=0.029$; Section~S5). In both
treatments the gain is carried by the day-of-year predictors; a correction given the model value
alone gains $+0.05$ under the instantaneous comparison and loses $-0.20$ under matched averaging
(Section~S5). What the correction learns from other monitors under the instantaneous comparison
is therefore mainly the mismatch in averaging, which all monitors share, and not something about
the atmosphere at the withheld site.

That no correction transfers under matched averaging is what Sections~\ref{sec:res-bias}
and~\ref{sec:res-subgrid} predict. What the reanalysis misses at a monitor is largely a local deficit;
other monitors do not carry that deficit, and the regional signal that they do carry is already
in the model.

\subsection{Consequences for assessing guideline exceedance}
\label{sec:res-exceed}

Of the 5,989 monitor-days evaluated, 5,046 (84\%) exceed the guideline in the observations. With
a base rate this high a product that always reported exceedance would be right on five days in
six, and only scores that discount chance agreement are informative (Table~\ref{tab:exceed}).

The reanalysis detects 3,722 of the 5,046 exceedances (POD 0.74, 95\% interval 0.64 to 0.83) and
reports exceedance on 487 of the 943 days that did not exceed (POFD 0.52). Its equitable threat
score is 0.09 (0.03 to 0.15) and its Heidke skill score 0.16 (0.06 to 0.25): above chance, but
small. Shifting the reanalysis by the monitor's own recent mean difference, which adds knowledge
of the local level but none of the timing, raises the equitable threat score to 0.28 (0.24 to
0.32). Persistence of the monitor's previous day reaches 0.39 (0.34 to 0.44). The ordering is the
one the previous sections lead to. The reanalysis carries part of the day-to-day signal, but
whether a given place exceeds the guideline on a given day depends mostly on a local level that
the reanalysis does not contain and that only a monitor at the place supplies.

\section{Discussion}

\subsection{Why column dust and surface PM$_{10}$ diverge}
\label{sec:disc-column}

The difference between a $0.75^{\circ}$ field and a point measurement contains model error and
representation error, and point data alone cannot fully separate them
\citep{schutgens2016,janjic2018}. The network nevertheless separates two parts of the error that
differ in cause. The first is shared by the coast and the interior and grows with the amount of
dust: the reanalysis compresses the day-to-day range, overestimating the cleanest days at Al~Ain
and underestimating the highest days everywhere, in the cool season as in the warm one. Because
desert dust dominates the regional aerosol burden \citep{gandham2020,shahid2026} and this part is
largest on the days with the most coarse aerosol over an independent sun photometer, it is an
error in the dust that the model brings to the surface of the region, not in the timing of dust
events, which the dust optical depth of the reanalysis follows closely.

Assimilation of optical depth constrains the total column and leaves the rest to the model. In the
CAMS system the analysis adjusts the total aerosol and distributes the increment over the species
in proportion to their share in the background \citep{inness2019}, and the model relates the mass
of each species and size bin to optical depth through prescribed optical properties
\citep{morcrette2009}. The comparison at the sun photometer places the regional deficit along this
chain (Section~S8). The first link holds: the total optical depth of the reanalysis is close to the
measured one (median ratio 1.07), as expected of a field that assimilates it. The second does not:
the reanalysis attributes to dust a smaller share of the column than the sun photometer attributes
to the coarse mode (0.31 against 0.48), so that its dust optical depth falls short of the measured
coarse-mode optical depth although the two rise and fall together. The third reduces the surface
value further: per unit of column dust the model carries 0.86 of the observed regional level. Both
links point the same way; medians over days do not show how they combine on a given day.

What sets the third link cannot be determined from these data. One contribution is built into
the product: the \pmten{} of the CAMS reanalysis is diagnosed from its aerosol bins, and of the
largest dust bin, with radii of 0.9 to 20~$\mu$m, only 40\% of the mass is counted as
\pmten{}\footnote{Copernicus Atmosphere Monitoring Service, answer on the ECMWF support forum,
26~September 2019, ``PM10 and PM25 global products'',
\url{https://forum.ecmwf.int/t/pm10-and-pm25-global-products/1084}; the formula given there for
the CAMS reanalysis counts 0.4 of the mass of the dust bin of 0.9 to 20~$\mu$m radius.}, so that a
given dust load yields less \pmten{} the more of it lies in that bin. How much of the deficit this
accounts for depends on the size distribution of the modelled dust near the surface, which was not
examined. Beyond the diagnostic, the ratio of surface level to
column dust mixes the vertical distribution of dust, that is the depth of the mixed layer and the
share of dust carried aloft, with the mass per unit optical depth, which depends on the size
distribution and is low for the large particles that carry much of the coarse mass. Evaluations of the
predecessor MACC reanalysis point to the vertical distribution. Against spaceborne lidar it
underestimated dust optical depth where loadings were high and, close to the Saharan and Middle
Eastern sources, underestimated dust extinction strongly between 1.2~km and about 3 to 5~km while
placing dust at heights up to 9~km \citep{georgoulias2018}; over the Sahel it reproduced the
variability of surface dust concentration but underestimated its daily and monthly means throughout
the year \citep{cuevas2015}, a pattern like the one found here. The vertical structure of reanalysis aerosol
products has also been evaluated against spaceborne lidar over the Tibetan Plateau
\citep{zhang2025tp}. Whether the current CAMS reanalysis distributes dust in the same way over the
Gulf was not tested: no profile measurement or size-resolved mass was available to separate these
effects, and as causes of the deficit all three remain hypotheses. The second link also has a caveat: the coarse mode of the sun photometer and the dust
of the reanalysis are not the same quantity, the first including coarse sea salt and the second
fine dust, which bias the comparison in opposite directions. For the evaluation of reanalysis
aerosol the result has a general consequence. Agreement in optical depth, and even agreement of
the dust optical depth in time, does not establish the surface concentration; the partition
between species and the relation of column to surface need to be checked against a surface
record.

\subsection{The local deficit of the conurbation}
\label{sec:disc-local}

The second part of the error is an error of place. What the reanalysis lacks is a set of offsets
that differ from one location to the next, are coherent over a few tens of kilometres, and vanish
on average at the inland city of Al~Ain. Over Abu Dhabi island and its suburbs four monitors in
two cells agree with each other and jointly exceed the model by more than 50~\ugm{}, so the
discrepancy cannot be assigned to the siting of one instrument. This deficit of the conurbation
relative to the interior, of about 50 to 60~\ugm{}, is
present already on the cleanest days and on the days with the least coarse aerosol, and is
not larger on the dustiest days. A deficit of this
kind is not explained by dust transport alone. It points to coarse particles of local origin, from traffic, construction, industry and resuspension \citep{thorpe2008,amato2009}, an urban
increment over the regional background \citep{lenschow2001} that a global model at this
resolution does not resolve, or to the composition of
the coastal cells, which extend over the sea as well as the city (Figure~\ref{fig:map}). The
explanations are not exclusive and the record cannot apportion the deficit between them. It does
weigh against wind-driven resuspension as the main cause, since the deficit does not grow with
the modelled wind speed or with dust, and against sea salt, which the reanalysis assigns a small
share of the optical depth; a wind field at $0.75^{\circ}$ does not, however, resolve the sea
breeze or the flow within the city.

Regulatory monitors are placed where exposure is expected to be high, so a network of this kind
overstates the deficit of a cell average and should not be read as an estimate of it. For the
question an assessment asks, namely what people at a given place breathe, the distinction matters
less than its size: between 31\% and 51\% of the spatial variance of \pmten{} in this
network is at scales the grid does not contain.

That the bias is spatially coherent but only loosely related to station type has a practical
corollary. A type-based adjustment, of the kind applied to make gridded products comparable with
urban or roadside monitors, would have been wrong here for at least one monitor in each
source-oriented class.

\subsection{The averaging interval of archived observations}

The values that an open archive holds for a regulatory network are whatever the network
publishes, and for particulate matter that may be the running mean on which an air-quality index
is based (Section~\ref{sec:running}), not the hourly concentration. In the present record the two are easy
to tell apart once the question is asked: hourly values that hardly change, increments that
reverse after 24~h, and gases from the same stations that behave otherwise. The same principle governs comparisons
between co-located instruments, where a systematic difference can arise entirely from what each
measured quantity includes \citep{you2026nai}. The check costs a
few lines of code and should precede any comparison of archived particulate data with a model.
The consequences of omitting it are not uniform. Mean bias and its spatial pattern were almost
unaffected here, because averaging does not change a mean. The correlation in time was
understated (0.37 instead of 0.50). And the evaluation of a statistical correction was
reversed: an apparent error reduction of 28--34\%, established at 16 and 17 of 19 monitors and
robust to every control applied, disappeared when model and observations were averaged alike.

\subsection{Why correction does not help, and what leave-one-station-out measures}

Statistical correction of chemical-transport output is well established where the corrected
product is used at the monitors that trained it \citep{djalalova2015,huang2017}, and the present
results are consistent with that use: knowledge of a monitor's own recent level raised the
exceedance score from 0.09 to 0.28. They do not support the further and common inference that a
correction fitted to a network can be mapped over areas without monitors. In this network the
leave-one-station-out gain was not established even before anything else was withheld (13 of 19
monitors improved), and it vanished once neighbours within 25~km were withheld or the training data
were separated from the target in time.
The requirement that a validation withhold what deployment will lack is familiar from
spatial prediction \citep{roberts2017,ploton2020,wadoux2021}, as is the inflation of skill by
near-duplicates of the training data under random splits in chemical machine learning
\citep{kim2026tas2r}, and the transferability of
statistical models between arid regimes has recently been examined for dust-related fine
particles \citep{hurynovich2026}. What the atmospheric setting adds is a physical reason for the
outcome. A correction can transfer only the part of the error that monitors share, and here the
shared part is the regional signal that the reanalysis already reproduces. The part that the
reanalysis misses is local, and no monitor elsewhere holds it. The conclusion applies to
corrections built, as this one is, from the model value, the date and the position; predictors
that describe the locality itself, such as land use, road density or emission inventories
\citep{zhang2018lur}, address the missing part directly and were not tested.

\subsection{Implications for assessment in arid, sparsely monitored regions}

Three implications follow for the use of reanalysis \pmten{} where monitors are few. First, a
reanalysis of this kind can indicate on which days \pmten{} over a region is high. It should not
be used to decide whether a particular place exceeded a guideline, and exceedance statistics
derived from it in a region where the guideline is exceeded on most days should be reported with
a chance-corrected score and the base rate. Second, the uncertainty attached to a grid value
used as a point estimate should include the sub-grid spread \citep{denby2011}, which in this network is as large
as the guideline value itself in the cells that contain strong local sources. Third, because the
missing component is a local level that changes slowly, a monitor informs its own site far more
than it informs a statistical model of the region; for a fixed number of instruments, coverage
of distinct source environments is worth more to assessment than replication within one. In
radiological surveys, likewise, an adaptive design guided by a statistical model of the field
found unvisited hot spots only where the field was correlated over distances larger than the
sampling spacing \citep{choi2027adaptive}.

\subsection{Limitations}
\label{sec:limits}

The study rests on one network in one emirate and on 19 months of record. The archive is thin in
spring and summer 2023, when it holds one value per day in three months, so the dust season is
represented mainly through the daily record. Days were classed by their regional \pmten{} level,
which follows the season. Dust was identified from a single sun photometer outside the network,
71 to 106~km from the monitors, whose Level~1.5 data have not received final calibration, and from
the weather reports of one airport inside it, which record dust on only 14 days, and the
wind and the aerosol composition were taken from the reanalysis itself, at a resolution that does
not resolve local circulations. No vertical profile of aerosol was available, so the relation between column
and surface could not be decomposed. Satellite dust optical depth, from MODIS for example, was not used to
check how well the single sun photometer represents the network, and \pmten{} was the only
particulate metric evaluated; fine particles (PM$_{2.5}$) were not. The archived values are integer-rounded 24-h running
means, so the study says nothing about diurnal variation or short peaks, and the alignment of the
averaging window, inferred from the data and from the national reporting convention, could not be
confirmed from station metadata; neither could the instruments in use during the period. Point monitors sited for regulatory purposes are not a random
sample of a cell. A single global product at $0.75^{\circ}$ was evaluated; a regional simulation at
kilometre scale with resolved local emissions would be the proper test of how much of the
sub-grid variance is recoverable by resolution \citep{ukhov2020}. The correction analysis has 19
units and one model family. The contrast between source-oriented and background station types was
formed after an omnibus test had failed and is reported as such. Finally, the exceedance analysis
classifies 24-h means in retrospect; it is not a verification of forecasts.

\section{Conclusions}

The \pmten{} values that the open archive holds for the 19 regulatory monitors of the Emirate of
Abu Dhabi are 24-h running means, and they exceed the WHO 24-h guideline value on most days at
every monitor. Compared over the same averaging interval and over 19 months, the CAMS reanalysis
followed the day-to-day course of the regional level (Spearman 0.67) and tracked single monitors
better than a distant monitor did, but it did not reproduce the differences between places: the
mean bias ranged from $-128$ to $+40$~\ugm{} and was organised in space. Its dust optical depth
followed an independent sun photometer in time and its total optical depth matched the measured
one, but it gave dust a smaller share of the column than the measured coarse mode and carried less
surface \pmten{} per unit of column dust, so that its regional deficit grew with dust, in both
seasons. Over the Abu Dhabi conurbation a local deficit of about 50 to 60~\ugm{} was added,
present already on the cleanest and on clear days and unrelated to modelled wind and dust, while the
inland city of Al~Ain was reproduced on average. Between 31\% and 51\% of the between-monitor variance of \pmten{} lay
inside single $0.75^{\circ}$ cells. A statistical correction built from the model value, the date and the position and trained on
other monitors gave no established gain at a monitor treated as unmonitored; the gain that appears when running means
are compared with instantaneous model values is produced by the mismatch in averaging. Judged by a chance-corrected score, the
reanalysis had little skill in separating days that exceeded the guideline at a monitor from days
that did not. In arid regions with strong local sources of coarse particles, a global reanalysis
describes when the regional background is high; statements about a particular place still require
a measurement there.
